% © 2025 Ing. David Jaroš — CC BY-NC-ND 4.0
%
% This work is licensed under a Creative Commons Attribution-NonCommercial-NoDerivatives
% 4.0 International License (CC BY-NC-ND 4.0).
%
% License History: Earlier drafts (up to v0.3) were released under CC BY 4.0.
% From v0.4 onward, all material is released under CC BY-NC-ND 4.0 to protect
% the integrity of the theoretical work during ongoing academic development.
%
% See LICENSE.md for full license text.

\ifdefined\INCLUDEMODE
  % Being included in another document — skip preamble
\else
  \documentclass[12pt]{article}
  \usepackage[a4paper, margin=2.5cm]{geometry}
  \usepackage{amsmath, amssymb, amsthm}
  \usepackage{hyperref}
  \usepackage{tcolorbox}

  \newtheorem{theorem}{Theorem}[section]
  \newtheorem{lemma}[theorem]{Lemma}
  \newtheorem{corollary}[theorem]{Corollary}
  \newtheorem{proposition}[theorem]{Proposition}
  \theoremstyle{definition}
  \newtheorem{definition}[theorem]{Definition}
  \theoremstyle{remark}
  \newtheorem{remark}[theorem]{Remark}

  \title{\textbf{Combined Variation vs.\ Termwise Separation\\
         in the UBT Action Principle}}
  \author{Ing.\ David Jaroš}
  \date{March 2026}

  \begin{document}
  \maketitle

  \begin{abstract}
  This note clarifies a subtle but important point in the variational derivation
  of Einstein's equations from the Unified Biquaternion Theory (UBT) action.
  When the metric $g_{\mu\nu}$ is promoted to a functional of the fundamental
  field $\Theta$ via $g_{\mu\nu} = g_{\mu\nu}[\Theta]$, the action becomes a
  functional of $\Theta$ alone, and its variation yields a single coupled
  Euler--Lagrange condition—\emph{not} two independently vanishing terms.
  Absence of automatic termwise separation is \emph{not} a mathematical flaw;
  it indicates that UBT may be broader than GR.  Einstein's equations are
  recovered on an admissible sector satisfying additional regularity and
  rank conditions, which are physically motivated and explicitly stated.
  UBT therefore \emph{contains} GR as a sector rather than contradicting it.
  \end{abstract}
\fi

% =========================================================================
\section{Combined Variation vs.\ Termwise Separation}
\label{sec:combined_vs_termwise}

\begin{tcolorbox}[colback=blue!5!white,colframe=blue!75!black,
                  title=Central Message of This Note]
The Euler--Lagrange equation obtained by varying the pure-$\Theta$ action
$\hat{S}[\Theta]$ is the \textbf{combined condition}
\begin{equation}
  \mathcal{E}_\Theta + J^{*}\mathcal{E}_g = 0,
  \tag{$\star$}
\end{equation}
where $\mathcal{E}_\Theta$ is the explicit-$\Theta$ Euler--Lagrange operator,
$\mathcal{E}_g$ is the metric Euler--Lagrange operator (the Einstein-sector
residual), and $J^{*}$ is the adjoint of the induced variation map
$J = \delta g / \delta\Theta$.

This is \emph{one} coupled equation, \emph{not} two equations $\mathcal{E}_\Theta = 0$
and $\mathcal{E}_g = 0$ imposed independently.  Termwise separation requires
additional structure (rank conditions, projections, or gauge choices) and is
not automatic.
\end{tcolorbox}

\subsection{Setup: Definitions}
\label{subsec:setup}

Let $\Theta: M^4 \to \mathcal{B}$ be the biquaternionic field, and let the
total action be
\begin{equation}
  S_{\text{total}}[g, \Theta]
  = S_{\text{EH}}[g] + S_\Theta[g, \Theta],
\end{equation}
where $S_{\text{EH}}$ is the Einstein--Hilbert action and $S_\Theta$ contains
all $\Theta$-dependent terms.

Define the operators:

\begin{definition}[Euler--Lagrange operators]
\label{def:el_operators}
\begin{align}
  \mathcal{E}_\Theta[\Theta, g]
  &:= \frac{\delta S_\Theta}{\delta\Theta}\bigg|_{g \text{ fixed}},
  \label{eq:E_Theta}\\[6pt]
  \mathcal{E}_g[g, \Theta]
  &:= \frac{\delta S_{\text{total}}}{\delta g^{\mu\nu}}\bigg|_{\Theta \text{ fixed}}
   = \frac{1}{16\pi G}\left(G_{\mu\nu} - 8\pi G\, T_{\mu\nu}[\Theta,g]\right),
  \label{eq:E_g}
\end{align}
and the \emph{induced variation map}
\begin{equation}
  J\,\delta\Theta := \frac{\delta g_{\mu\nu}[\Theta]}{\delta\Theta}\,\delta\Theta
  \qquad\text{(a tensor-valued linear map)}.
  \label{eq:J_map}
\end{equation}
The formal $L^2$-adjoint of $J$ is denoted $J^*$.
\end{definition}

% =========================================================================
\section{Why This Is Not Separation of Variables or Separation of Constants}
\label{sec:not_separation}

\subsection{The Pure-$\Theta$ Functional and Its Variation}
\label{subsec:pure_theta}

Substituting the emergent metric $g_{\mu\nu}[\Theta]$ into
$S_{\text{total}}$, we obtain the pure-$\Theta$ functional
\begin{equation}
  \hat{S}[\Theta]
  := S_{\text{total}}\bigl[g[\Theta],\, \Theta\bigr].
  \label{eq:S_hat}
\end{equation}

Varying with respect to $\Theta$ and applying the chain rule:
\begin{equation}
  \frac{\delta\hat{S}}{\delta\Theta}\,\delta\Theta
  = \underbrace{\frac{\delta S_{\text{total}}}{\delta\Theta}\bigg|_g}_{\mathcal{E}_\Theta}
    \delta\Theta
  + \underbrace{\frac{\delta S_{\text{total}}}{\delta g^{\mu\nu}}}_{\mathcal{E}_g^{\mu\nu}}
    \underbrace{\frac{\delta g^{\mu\nu}[\Theta]}{\delta\Theta}}_{J}\delta\Theta.
  \label{eq:variation_chain}
\end{equation}

The stationarity condition $\delta\hat{S} = 0$ for all admissible variations
$\delta\Theta$ therefore reads:
\begin{equation}
  \boxed{
    \mathcal{E}_\Theta + J^{*}\mathcal{E}_g = 0.
  }
  \label{eq:combined_EL}
\end{equation}

\begin{remark}[Why separation does not follow]
\label{rem:no_separation}
Even though Eq.~\eqref{eq:combined_EL} holds for all $\delta\Theta$, one
\emph{cannot} conclude that $\mathcal{E}_\Theta = 0$ and $\mathcal{E}_g = 0$
independently.  The correct statement is only that the \emph{linear combination}
$\mathcal{E}_\Theta + J^{*}\mathcal{E}_g$ vanishes.  This is analogous to a
linear equation $a\,x + b\,y = 0$: knowing this holds for all vectors
$(x, y)$ would allow separation, but here $\delta\Theta$ is a single variation
and the two terms are coupled.
\end{remark}

\begin{remark}[No resemblance to separation of constants]
This structure should \emph{not} be described as ``separation of variables''
or ``separation of constants'' in any sense.  Those terms refer to specific
algebraic decoupling techniques that have no bearing here.
The correct description is: the combined variational condition is a
\emph{single first-order functional differential equation} in $\Theta$.
\end{remark}

\subsection{When Does $\mathcal{E}_g = 0$ Follow Separately?}
\label{subsec:when_sep}

Separate vanishing $\mathcal{E}_g = 0$ (i.e.\ Einstein's equations
$G_{\mu\nu} = 8\pi G\, T_{\mu\nu}$) follows from Eq.~\eqref{eq:combined_EL}
under any one of the following additional conditions:

\begin{enumerate}
  \item \textbf{Rank condition on $J$:} If the map $J$ is locally surjective
    (full column rank on the relevant function space) and $\mathcal{E}_\Theta$
    vanishes on the solution, then $J^{*}\mathcal{E}_g = 0$ together with
    surjectivity implies $\mathcal{E}_g = 0$.

  \item \textbf{Gauge-reduced injectivity:} If, after quotienting by gauge
    redundancies, the map $\Theta \mapsto g[\Theta]$ is locally invertible
    (non-degenerate Jacobian on the physical sector), then variations of
    $g$ induced by $\delta\Theta$ span the full metric variation space, and
    independent metric stationarity follows.

  \item \textbf{Projection to metric observables:} If one restricts attention
    to $\Theta$ configurations whose only observable degrees of freedom are
    metric (all additional $\Theta$ modes are either gauge or frozen), then
    the contracted form of Eq.~\eqref{eq:combined_EL} reduces to
    $\mathcal{E}_g = 0$.

  \item \textbf{On-shell $\mathcal{E}_\Theta = 0$:} If $\Theta$ satisfies its
    own field equation separately (e.g.\ on a constrained sub-sector or after
    solving for the matter equations of motion), then Eq.~\eqref{eq:combined_EL}
    reduces to $J^{*}\mathcal{E}_g = 0$, and the rank condition above suffices.
\end{enumerate}

\begin{tcolorbox}[colback=yellow!10!white,colframe=orange!80!black,
                  title=Warning: Do Not Assume Arbitrary Independent Variation]
It is mathematically \emph{incorrect} to argue:
\begin{quote}
``Since the variation holds for arbitrary $\delta\Theta$, each term in
Eq.~\eqref{eq:combined_EL} must vanish separately.''
\end{quote}
The two terms $\mathcal{E}_\Theta\,\delta\Theta$ and
$J^{*}\mathcal{E}_g\,\delta\Theta$ are \emph{not} independently varied;
they both depend on the same variation $\delta\Theta$.  Separate vanishing
would require $\delta\Theta$ to vary each term independently, which is
generally impossible.
\end{tcolorbox}

% =========================================================================
\section{The Fundamental Equation $\mathcal{E}_\Theta + J^*\mathcal{E}_g = 0$}
\label{sec:fundamental_equation}

Equation~\eqref{eq:combined_EL} is the \textbf{fundamental variational
statement} of UBT in the $\Theta$-only formulation.  Its properties:

\begin{enumerate}
  \item It is a \emph{necessary} condition for stationarity of $\hat{S}[\Theta]$.
  \item It encodes both the $\Theta$-field dynamics and the Einstein-sector
    constraint in a single equation.
  \item It reduces to \emph{Einstein's equations plus the $\Theta$ field
    equation} when the rank condition holds (Section~\ref{subsec:when_sep}).
  \item It is \emph{weaker} than the system
    $\{\mathcal{E}_\Theta = 0,\; \mathcal{E}_g = 0\}$:
    the combined condition admits additional solutions not present in
    Einstein gravity.  These additional solutions live in the
    \emph{non-GR sector} of UBT.
\end{enumerate}

% =========================================================================
\section{When Separate Recovery of $\mathcal{E}_g = 0$ Does and Does Not Follow}
\label{sec:recovery}

\subsection{Does Follow}

\begin{theorem}[GR recovery on admissible metric-generating sector]
\label{thm:gr_recovery}
Suppose $\Theta$ lies in the admissible class $\mathcal{A}_{\mathrm{UBT}}$ defined
by the conditions in \textup{Section~\ref{subsec:when_sep}}.  Then the combined
Euler--Lagrange condition~\eqref{eq:combined_EL} implies
\begin{equation}
  G_{\mu\nu} = 8\pi G\, T_{\mu\nu}[\Theta, g[\Theta]],
  \label{eq:einstein}
\end{equation}
i.e.\ Einstein's field equations hold on this sector.
\end{theorem}

\begin{proof}[Proof sketch]
Under the rank condition (Case 1 or 2 of Section~\ref{subsec:when_sep}), any
variation $\delta g^{\mu\nu}$ can be written as $J\,\delta\Theta$ for some
$\delta\Theta$.  Therefore $\langle \mathcal{E}_g,\, J\,\delta\Theta \rangle = 0$
for all $\delta\Theta$ implies $\langle J^*\mathcal{E}_g,\, \delta\Theta \rangle = 0$
for all $\delta\Theta$.  Surjectivity of $J$ then gives $\mathcal{E}_g = 0$. \qed
\end{proof}

\subsection{Does Not Follow (General Case)}

In the general case, outside $\mathcal{A}_{\mathrm{UBT}}$:
\begin{itemize}
  \item The kernel of $J$ may be non-trivial: some $\delta\Theta$ directions
    do not change the metric.  These ``metric-invisible'' variations are
    genuine UBT degrees of freedom absent from GR.
  \item The combined condition may be satisfied by configurations where
    $\mathcal{E}_g \neq 0$ (non-GR solutions of UBT).
  \item This is \emph{not a defect}: it signals that UBT contains GR as a
    proper subsector while admitting additional dynamics.
\end{itemize}

% =========================================================================
\section{Three Notions of GR Recovery}
\label{sec:three_notions}

Depending on the strength of the conditions imposed, three distinct levels of
GR recovery can be distinguished.  These are elaborated in
\texttt{reports/gr\_recovery\_levels.md}; a brief summary follows.

\begin{enumerate}
  \item \textbf{Exact projected recovery:} A mathematically clean
    projection/contraction theorem maps the $\Theta$ equation to
    $G_{\mu\nu} = 8\pi G\, T_{\mu\nu}$ exactly.
    \emph{Status: open problem (requires full rank condition).}

  \item \textbf{Constrained sector recovery:} On the admissible class
    $\mathcal{A}_{\mathrm{UBT}}$ (metric-generating sector with gauge
    reduction), the induced dynamics reproduces Einstein gravity.
    \emph{Status: the current best-established description;
    conditions are stated explicitly in
    \texttt{GR\_closure/gr\_sector\_conditions.tex}.}

  \item \textbf{Low-energy effective limit recovery:} After freezing
    additional $\Theta$ degrees of freedom that decouple at macroscopic
    scales, the theory reduces to GR in the infrared.
    \emph{Status: plausible; requires analysis of the $\Theta$-mode
    spectrum.}
\end{enumerate}

% =========================================================================
\section{Why UBT $\supseteq$ GR Is Not a Contradiction}
\label{sec:no_contradiction}

\begin{tcolorbox}[colback=green!8!white,colframe=green!60!black,
                  title=Consistency Statement]
There is \textbf{no contradiction} between UBT and GR.  UBT can reproduce GR
on a physical sector while still containing additional non-GR degrees of
freedom.  The absence of automatic termwise separation is a feature, not a
bug: it indicates that UBT is \emph{richer} than GR, not inconsistent with it.
\end{tcolorbox}

The logic is standard in fundamental physics:

\begin{itemize}
  \item \textbf{Analogy 1 — Tetrad vs.\ metric variation:} In first-order
    (Palatini or tetrad) formulations of GR, varying the tetrad $e^a{}_\mu$
    and the connection $\omega^{ab}{}_\mu$ independently yields Einstein's
    equations plus a torsion-free condition.  Neither condition follows from
    the other without additional structure; yet the combined system is fully
    consistent.  The UBT situation is analogous: $\Theta$ is the fundamental
    object, and the metric is a derived (induced) quantity.

  \item \textbf{Analogy 2 — Constrained field redefinitions in effective field
    theory (EFT):} In EFT, integrating out heavy fields yields a low-energy
    effective action whose field equations are not termwise equal to those of
    the UV theory; they agree only on-shell or in a specific limit.  GR
    recovery from UBT is structurally of this type.

  \item \textbf{General principle:} A deeper theory is expected to contain
    the standard theory as a sector or limit, not to reproduce it in every
    configuration.  String theory contains GR as a low-energy limit;
    loop quantum gravity recovers GR in the classical limit; UBT recovers
    GR on the metric-generating sector.  None of these is a contradiction.
\end{itemize}

\bigskip
\noindent
\textbf{Correct language:}
\begin{itemize}
  \item[\checkmark] ``UBT generalises GR'' — accurate.
  \item[\checkmark] ``GR is recovered as a sector/limit/projection of UBT'' — accurate.
  \item[\checkmark] ``UBT extends GR with additional degrees of freedom'' — accurate.
  \item[$\times$] ``UBT contradicts GR'' — incorrect.
  \item[$\times$] ``Absence of termwise separation weakens UBT fatally'' — incorrect.
  \item[$\times$] ``If exact GR is not separately extracted, UBT fails'' — incorrect.
\end{itemize}

% =========================================================================

\ifdefined\INCLUDEMODE
  % no end{document} when included
\else
  \end{document}
\fi
